\documentclass[12pt]{article}
\usepackage[utf8]{inputenc}
\usepackage{amsmath, amssymb}
\usepackage{enumitem}
\usepackage[spanish]{babel}

\title{Examen de Física -- Nivel Universitario}
\author{}
\date{}

\begin{document}

\maketitle

\noindent
\textbf{Instrucciones:} Selecciona la opción correcta en cada pregunta. Las respuestas numéricas difieren sólo en el último decimal; las conceptuales presentan matices sutiles pero determinantes.

\begin{enumerate}

\item Un bloque de $2\ \text{kg}$ se mueve sobre una superficie horizontal sin fricción con velocidad $4\ \text{m/s}$. Súbitamente se le aplica una fuerza constante de $5\ \text{N}$ en dirección opuesta al movimiento durante $3\ \text{s}$. ¿Cuál es la magnitud de su velocidad final?
\begin{enumerate}[label=(\alph*)]
\item $3.5\ \text{m/s}$
\item $3.0\ \text{m/s}$
\item $2.5\ \text{m/s}$
\item $4.0\ \text{m/s}$
\end{enumerate}

\item Un proyectil es lanzado con rapidez inicial de $30\ \text{m/s}$ y ángulo de $53^\circ$ sobre la horizontal. Despreciando la resistencia del aire, ¿cuál es la altura máxima alcanzada? ($g = 9.8\ \text{m/s}^2$, $\sin53^\circ=0.8$)
\begin{enumerate}[label=(\alph*)]
\item $28.8\ \text{m}$
\item $29.4\ \text{m}$
\item $27.5\ \text{m}$
\item $30.1\ \text{m}$
\end{enumerate}

\item Una partícula realiza un movimiento circular uniforme con radio $0.5\ \text{m}$ y frecuencia $2\ \text{Hz}$. ¿Cuál es la magnitud de su aceleración centrípeta?
\begin{enumerate}[label=(\alph*)]
\item $78.9\ \text{m/s}^2$
\item $79.2\ \text{m/s}^2$
\item $78.5\ \text{m/s}^2$
\item $80.0\ \text{m/s}^2$
\end{enumerate}

\item Dos cuerpos de masas $3\ \text{kg}$ y $5\ \text{kg}$ están unidos por una cuerda ligera que pasa por una polea ideal. Si se dejan en libertad, ¿cuál es la tensión en la cuerda? ($g = 9.8\ \text{m/s}^2$)
\begin{enumerate}[label=(\alph*)]
\item $36.8\ \text{N}$
\item $37.5\ \text{N}$
\item $36.0\ \text{N}$
\item $37.0\ \text{N}$
\end{enumerate}

\item Un auto de $1200\ \text{kg}$ viaja a $20\ \text{m/s}$ y frena con una fuerza constante de $3000\ \text{N}$. ¿Qué distancia recorre hasta detenerse?
\begin{enumerate}[label=(\alph*)]
\item $80.0\ \text{m}$
\item $79.8\ \text{m}$
\item $80.2\ \text{m}$
\item $78.5\ \text{m}$
\end{enumerate}

\item Se tiene un péndulo simple de $1\ \text{m}$ de longitud. ¿Cuál es su periodo en la Luna ($g_{\text{luna}} = 1.6\ \text{m/s}^2$)?
\begin{enumerate}[label=(\alph*)]
\item $4.96\ \text{s}$
\item $5.02\ \text{s}$
\item $4.85\ \text{s}$
\item $5.10\ \text{s}$
\end{enumerate}

\item Un disco de momento de inercia $0.2\ \text{kg·m}^2$ gira a $10\ \text{rad/s}$. Se le aplica un torque constante de $0.5\ \text{N·m}$ durante $4\ \text{s}$. ¿Cuál es su velocidad angular final?
\begin{enumerate}[label=(\alph*)]
\item $20.0\ \text{rad/s}$
\item $19.5\ \text{rad/s}$
\item $20.5\ \text{rad/s}$
\item $18.0\ \text{rad/s}$
\end{enumerate}

\item Un gas ideal experimenta una expansión isotérmica a $300\ \text{K}$ desde $2\ \text{L}$ hasta $5\ \text{L}$. Si la presión inicial era $4\ \text{atm}$, ¿cuál es el trabajo realizado por el gas? ($1\ \text{atm·L} = 101.3\ \text{J}$)
\begin{enumerate}[label=(\alph*)]
\item $743\ \text{J}$
\item $745\ \text{J}$
\item $740\ \text{J}$
\item $748\ \text{J}$
\end{enumerate}

\item Un recipiente de $3\ \text{m}^3$ contiene $2\ \text{mol}$ de gas ideal a $27^\circ\text{C}$. Si se calienta a presión constante hasta $127^\circ\text{C}$, ¿cuál es el volumen final?
\begin{enumerate}[label=(\alph*)]
\item $4.00\ \text{m}^3$
\item $3.98\ \text{m}^3$
\item $4.02\ \text{m}^3$
\item $3.95\ \text{m}^3$
\end{enumerate}

\item Un fluido incompresible fluye por una tubería horizontal que se estrecha de $4\ \text{cm}^2$ a $2\ \text{cm}^2$. Si la velocidad en la sección ancha es $1\ \text{m/s}$ y la presión $200\ \text{kPa}$, ¿cuál es la presión en la sección estrecha? (densidad $= 1000\ \text{kg/m}^3$)
\begin{enumerate}[label=(\alph*)]
\item $198.5\ \text{kPa}$
\item $199.0\ \text{kPa}$
\item $200.5\ \text{kPa}$
\item $197.5\ \text{kPa}$
\end{enumerate}

\item Un objeto de $0.5\ \text{kg}$ flota en agua (densidad $1000\ \text{kg/m}^3$) con el $80\%$ de su volumen sumergido. ¿Cuál es la densidad del objeto?
\begin{enumerate}[label=(\alph*)]
\item $800.0\ \text{kg/m}^3$
\item $802.0\ \text{kg/m}^3$
\item $798.5\ \text{kg/m}^3$
\item $801.0\ \text{kg/m}^3$
\end{enumerate}

\item Una carga puntual de $+2\ \mu\text{C}$ se coloca en el origen. ¿Cuál es el potencial eléctrico a $3\ \text{cm}$ de distancia? ($k = 9\times10^9\ \text{N·m}^2/\text{C}^2$)
\begin{enumerate}[label=(\alph*)]
\item $600\ \text{kV}$
\item $599\ \text{kV}$
\item $601\ \text{kV}$
\item $598\ \text{kV}$
\end{enumerate}

\item Un capacitor de $10\ \mu\text{F}$ se conecta a una batería de $12\ \text{V}$. Luego se desconecta y se conecta a otro capacitor descargado de $5\ \mu\text{F}$. ¿Cuál es la tensión final común?
\begin{enumerate}[label=(\alph*)]
\item $8.0\ \text{V}$
\item $8.1\ \text{V}$
\item $7.9\ \text{V}$
\item $8.2\ \text{V}$
\end{enumerate}

\item Una espira circular de radio $0.1\ \text{m}$ está en un campo magnético uniforme de $0.5\ \text{T}$ perpendicular a su plano. Si el campo se reduce a cero en $0.2\ \text{s}$, ¿cuál es la fem inducida promedio?
\begin{enumerate}[label=(\alph*)]
\item $78.5\ \text{mV}$
\item $79.0\ \text{mV}$
\item $78.0\ \text{mV}$
\item $80.0\ \text{mV}$
\end{enumerate}

\item Un protón entra con velocidad $2\times10^6\ \text{m/s}$ perpendicularmente a un campo magnético de $0.3\ \text{T}$. ¿Cuál es el radio de su trayectoria? (masa protón $= 1.67\times10^{-27}\ \text{kg}$, carga $= 1.6\times10^{-19}\ \text{C}$)
\begin{enumerate}[label=(\alph*)]
\item $6.96\ \text{cm}$
\item $6.90\ \text{cm}$
\item $7.02\ \text{cm}$
\item $6.85\ \text{cm}$
\end{enumerate}

\item Un circuito $RL$ serie tiene $R = 10\ \Omega$ y $L = 2\ \text{H}$. ¿Cuál es la constante de tiempo del circuito?
\begin{enumerate}[label=(\alph*)]
\item $0.20\ \text{s}$
\item $0.19\ \text{s}$
\item $0.21\ \text{s}$
\item $0.18\ \text{s}$
\end{enumerate}

\item Una onda armónica tiene ecuación $y(x,t) = 0.05\sin(4\pi x - 20\pi t)$ (SI). ¿Cuál es su velocidad de propagación?
\begin{enumerate}[label=(\alph*)]
\item $5.0\ \text{m/s}$
\item $5.1\ \text{m/s}$
\item $4.9\ \text{m/s}$
\item $5.2\ \text{m/s}$
\end{enumerate}

\item Dos fuentes sonoras emiten ondas coherentes con frecuencia $680\ \text{Hz}$. Si están separadas $0.5\ \text{m}$ y una persona se ubica a $2\ \text{m}$ en línea recta perpendicular al centro, ¿cuál es la diferencia de caminos aproximada? (velocidad del sonido $340\ \text{m/s}$)
\begin{enumerate}[label=(\alph*)]
\item $0.062\ \text{m}$
\item $0.060\ \text{m}$
\item $0.065\ \text{m}$
\item $0.058\ \text{m}$
\end{enumerate}

\item Un objeto de $2\ \text{cm}$ de altura se coloca a $15\ \text{cm}$ de una lente convergente de distancia focal $10\ \text{cm}$. ¿Cuál es la altura de la imagen?
\begin{enumerate}[label=(\alph*)]
\item $4.0\ \text{cm}$
\item $3.9\ \text{cm}$
\item $4.1\ \text{cm}$
\item $3.8\ \text{cm}$
\end{enumerate}

\item La luz roja de longitud de onda $650\ \text{nm}$ en el aire entra en un vidrio con índice de refracción $1.5$. ¿Cuál es la longitud de onda en el vidrio?
\begin{enumerate}[label=(\alph*)]
\item $433.3\ \text{nm}$
\item $434.0\ \text{nm}$
\item $432.5\ \text{nm}$
\item $435.0\ \text{nm}$
\end{enumerate}

\item La función trabajo del potasio es $2.3\ \text{eV}$. ¿Cuál es la frecuencia umbral para el efecto fotoeléctrico? ($h = 4.14\times10^{-15}\ \text{eV·s}$)
\begin{enumerate}[label=(\alph*)]
\item $5.56\times10^{14}\ \text{Hz}$
\item $5.55\times10^{14}\ \text{Hz}$
\item $5.57\times10^{14}\ \text{Hz}$
\item $5.54\times10^{14}\ \text{Hz}$
\end{enumerate}

\item Un electrón se acelera desde el reposo mediante una diferencia de potencial de $1000\ \text{V}$. ¿Cuál es su longitud de onda de De Broglie? ($h = 6.63\times10^{-34}\ \text{J·s}$, $m_e = 9.11\times10^{-31}\ \text{kg}$, $e = 1.6\times10^{-19}\ \text{C}$)
\begin{enumerate}[label=(\alph*)]
\item $3.88\times10^{-11}\ \text{m}$
\item $3.87\times10^{-11}\ \text{m}$
\item $3.89\times10^{-11}\ \text{m}$
\item $3.86\times10^{-11}\ \text{m}$
\end{enumerate}

\item Un núcleo de uranio-$238$ emite una partícula alfa. ¿Cuál es el número atómico y másico del núcleo resultante?
\begin{enumerate}[label=(\alph*)]
\item $Z=90$, $A=234$
\item $Z=91$, $A=234$
\item $Z=90$, $A=235$
\item $Z=92$, $A=234$
\end{enumerate}

\item Un satélite orbita la Tierra en una órbita circular. Si su rapidez orbital se duplica (sin cambiar la dirección), ¿qué ocurre con la órbita?
\begin{enumerate}[label=(\alph*)]
\item Pasa a una órbita elíptica con el perigeo en el punto de cambio.
\item Se escapa a una trayectoria parabólica.
\item Pasa a una órbita elíptica con el apogeo en el punto de cambio.
\item Continúa en órbita circular de menor radio.
\end{enumerate}

\item Un imán se acerca a una espira conductora cerrada. ¿Cuál de las siguientes afirmaciones es \textbf{correcta} según la ley de Lenz?
\begin{enumerate}[label=(\alph*)]
\item La corriente inducida crea un campo magnético que se opone al aumento del flujo, pero si el imán se aleja, la corriente ayuda a disminuir el flujo.
\item La corriente inducida siempre crea un campo que se opone al movimiento del imán, independientemente de si se acerca o se aleja.
\item La corriente inducida tiende a mantener constante el flujo magnético, por lo que si el imán se acerca, la corriente atrae al imán.
\item La corriente inducida crea un campo que siempre tiene la misma dirección que el campo del imán, para reforzarlo.
\end{enumerate}

\item Dos resortes idénticos se conectan en serie y luego en paralelo a una misma masa. ¿Cómo se comparan los periodos de oscilación?
\begin{enumerate}[label=(\alph*)]
\item El periodo en serie es el doble que en paralelo.
\item El periodo en serie es la mitad que en paralelo.
\item El periodo en serie es $\sqrt{2}$ veces el periodo en paralelo.
\item El periodo en serie es $2$ veces el periodo en paralelo.
\end{enumerate}

\item Una onda electromagnética plana se propaga en el vacío. Si el campo eléctrico oscila en la dirección $y$, ¿en qué dirección oscila el campo magnético?
\begin{enumerate}[label=(\alph*)]
\item En la misma dirección $y$, pero en fase.
\item En la dirección $z$, perpendicular a la propagación y a $\mathbf{E}$.
\item En la dirección $x$ (dirección de propagación).
\item En una dirección que forma $45^\circ$ con el plano $xy$.
\end{enumerate}

\end{enumerate}

\end{document}